\scanpage{27}\mspage{21}
% printed form: music manuscript paper, with blank staves behind the handwriting.
% @ upper staff, continuation of c | text | Ψ tilde de Ψ
$\widetilde\Psi(V,t)$ de $\Psi(V,t)=\operatorname{Ker}\bigl(V^t\overset{\varepsilon}{\longrightarrow}V\bigr)$ par $\mathbb F_2$ \inserted{$(x_i)_{i\in t}\longmapsto\sum x_i$} -- extension à \unsure{splittings|splittages|splitting} \inserted{précisés} sur les $N_i=\operatorname{Ker}\ \Psi(V,t)\overset{p_i}{\longrightarrow}V$ $(i\in t)$ i.e. on a des sous-groupes $N'_i\subset\widetilde\Psi(V,t)$, $N'_i\longrightarrow N_i$. Rappelons que $\widetilde\Psi(V,t)=\bigwedge_{i\in t}p_i^*(\widetilde V)$, où $\widetilde V$ est l'ext. de $V$ par $\mathbb F_2$ qui ne splitte sur aucun sous-groupe $\ne0$ de $V$. On a, alors \unsure{pr tt|notant|pour} $i\in t$

% @ upper middle, eight and four element sets | display | P̃ᵢ, Pᵢ
\[
\begin{aligned}
\overset{\text{``8''}}{\widetilde P_i}&=\widetilde P/N'_i\qquad i\in t\\
P_i&=\widetilde P/\widetilde N_i
\end{aligned}
\]
\margin{$E(i)^*$ ``}\margin{\deleted{\unsure{$E(i)$}}}
% @ right of quotients | text | où Ñᵢ
(où $\widetilde N_i=$ image inverse de $N_i$ dans $\widetilde\Psi$) $=N'_i.\{1,\sigma\}$, \unsure{$\sigma$} -- élém. \unsure{val.} \inserted{$\ne1$ du noyau de $\widetilde\Psi\longrightarrow\Psi$})
% @ centre, set of vertices | display | C ≃ t ⨿
\[
C\simeq t\amalg\left(\coprod_{i\in t}\widetilde P_i\right)
\]

% @ centre, heading | text | La loi du graphe
La loi du graphe est définie ainsi

% @ centre, a | text | les i ∈ t
a) les $i\in t$ sont deux à deux liés

% @ centre, b | text | i est lié
b) $i$ est lié aux él. de $\widetilde P_i$, et non à ceux de $\widetilde P_j$ si $j\ne i$

% @ lower middle, c | text | Si x, y
c) Si \deleted{$x\in\widetilde P_i$,} $\tilde x,\tilde y\in\widetilde P_i$, $x$ et $y$ sont liés sss $y=\sigma x$

% @ lower middle, d | text | Si x̃ ∈ P̃ᵢ
d) Si $\tilde x\in\widetilde P_i,\tilde y\in\widetilde P_j$, $i\ne j$, alors $x$ et $y$ sont liés sss $\exists\,\tilde{\unsure{\vartheta|\theta|\gamma}}\in\widetilde P$ tel que \deleted{$p_i(\tilde\vartheta)=x$,} $p_i(\tilde\vartheta)=\tilde x$, $p_j(\tilde\vartheta)=\tilde y$.

% @ left margin, first slanted note | text | NB P̃ ≃ ens des triangles
\margin{\uline{NB} $\widetilde P\simeq$ ens des triangles $t'$ disjoints de $t$}% slanted
% @ left margin, second slanted note | text | donc Coro Se donner
\margin{donc \uline{Coro} Se donner $(C,t,t')$ avec \unsure{$(t,t')$} triangles disjoints, équivaut à se donner $(t,V)$ ($t$ ens de card 3, $V$ vectoriel de dim 2 sur $\mathbb F_2$)}% slanted

% @ foot, corollary | text | Le groupe des autom
\uline{Coro} Le groupe des autom de $(C,t)$ \deleted{est} est isom au groupe des \deleted{couples} \inserted{triplets} $(\tau,\rho,\varphi)$ d'un autom $\rho$ de $\unsure{\tau|t}$ de $t$, d'un autom $\rho$ de $V$, et d'un \deleted{\unsure{$(\tau,\rho)$}} $\lambda_{\tau,\rho}$-autom de $\unsure{\widehat P|\widetilde P}$, où $\lambda_{\tau,\rho}$ est l'autom de \deleted{\unsure{$V_t$|$\Psi$}} $\Psi(V,t)$ induit par \deleted{$(\rho,\tau)$} $(\tau,\rho)$. Donc
% @ bottom staff, exact sequence | diagram | 1, P̃, W, S_t × S_V*, 1
\diagram{\begin{tikzcd}[ampersand replacement=\&,column sep=small]
1\arrow[r]\&\widetilde{\unsure{P|\Gamma}}\arrow[r]\&W_t\arrow[r]\&\overset{\simeq\mathfrak S_3\times\mathfrak S_3}{\unsure{\mathfrak S|\mathscr S}_t\times\mathfrak S_{V^*}}\arrow[r,dash]\&1
\end{tikzcd}}
% archive sheet number: 26
